\documentclass[11pt]{article}
\usepackage{fontspec,amsmath,amssymb,amsthm,hyperref}
\usepackage[margin=1in]{geometry}
\setlength{\emergencystretch}{3em}
\title{The KLMC law bound with the genuine initial Wasserstein distance}
\author{LangevinLean mathematical exposition}
\date{}
\begin{document}\maketitle
\noindent Source module: \texttt{KineticLangevin/CanonicalKLMCBound.lean}. Compile this proof document with XeLaTeX or LuaLaTeX. The accompanying HTML pairs the mathematical argument with the exact Lean statements and records the scope of the formalization.


The numerical law is the pushforward of the implemented KLMC chain on the actual independent Gaussian-velocity/Brownian product source. Every initial transport coupling yields this same output law. The complete constructed coupling proof gives the paper's first-order bound, and taking the actual transport infimum replaces the supplied coupling cost by the true initial Wasserstein distance. This intermediate module makes the sole stationary global-reference law property explicit; the final Gibbs theorem discharges it from the actual Brownian source.


\section{Fix the actual numerical output law}


The definition uses the original pointwise KLMC method, the standard Gaussian initial velocity and the exact canonical correlated Brownian innovations. The numerical L\ensuremath{{}^{2}} chain is identified with the generic product-law construction by equality of the actual innovation variables. The established output-law theorem then proves that its pushforward is independent of which initial position coupling was chosen.



\noindent Lean declarations: \texttt{klmcLaw}, \texttt{numericalChain\_outputLaw}.

\section{Transfer one actual global stationary phase law to the reference segments}


The named stationary property concerns the full constructed kinetic flow: at every time its phase law is \(\mathcal N(0,I_d)\otimes\pi\). It contains no numerical convergence conclusion. Exact chain-to-global identification transfers its position marginal to every reference grid point. Exact segment-to-global identification transfers its Gaussian velocity marginal to every intermediate time in a reference segment. Lean obtains these marginals by composition with the product-measure projections.



\noindent Lean declarations: \texttt{GlobalReferenceStationary}, \texttt{referenceChain\_position\_law}, \texttt{referenceSegment\_velocity\_law}.

\section{Apply the proved coupling estimate to one actual initial plan}


The initial position L\ensuremath{{}^{2}} discrepancy equals the true cost of the supplied transport coupling; this follows from the concrete extension's representative and norm identities. The segment specification furnishes all analytic local-error premises internally. Together with the derived stationary position and Gaussian velocity marginals, the quantitative L\ensuremath{{}^{2}} KLMC theorem gives \(W_2(\nu_n,\pi)\le\sqrt2(1-3mh/(4\gamma))^n\operatorname{cost}(c)+Mh\sqrt{2d}/m\). Output-law equality and initial-cost equality rewrite this into a bound for the actual numerical law.



\noindent Lean declarations: \texttt{initial\_position\_cost}, \texttt{klmc\_bound\_of\_globalStationarity}.

\section{Prove output finiteness and take the transport infimum}


Finite initial Wasserstein distance supplies an actual finite-cost initial coupling. Its numerical and reference states are genuine L\ensuremath{{}^{2}} variables, so their distance gives a finite transport bound on the output law. This proves output finiteness before real-valued transport arithmetic is used. The scalar recurrence estimate proves the contraction factor is nonnegative under the paper's step restrictions. Lean then applies the proved infimum theorem over all finite-cost couplings, obtaining the same estimate with \(W_2(\mu,\pi)\) in place of the cost of a chosen plan. No existence of an optimal coupling is assumed.



\noindent Lean declarations: \texttt{klmcLaw\_wasserstein2\_ne\_top}, \texttt{klmc\_wasserstein\_bound\_of\_globalStationarity}.

\end{document}
